\begin{definition}[Viscosity sub-solution, Definition~2.1]
Let $\mathbb G$ be a Carnot setting (\noderef{CarnotSetting}) with horizontal frame
$\sigma$, let $A:\mathbb R^m\to\mathbb R^{m\times m}$ and $H:\mathbb R^m\to\mathbb R$, let
$\Omega\subset\mathbb R^n$ and let $u:\mathbb R^n\to\mathbb R$. We say that $u$ is a
\emph{viscosity sub-solution of $\mathcal Lu=0$ in $\Omega$}, written
$\mathrm{IsViscSubsol}\ \mathbb G\ A\ H\ \Omega\ u$, if
\begin{enumerate}
\item $u$ is upper semi-continuous on $\Omega$; and
\item for every $x\in\Omega$ and every $\varphi\in C^{2}(\Omega)$ with
  $x\in\mathrm{argmax}_\Omega(u-\varphi)$, every second-order jet $(D,B)$ of $\varphi$ at
  $x$ (\noderef{TwiceDiffAt}) with $\mathrm{hgrad}\ \sigma\ x\ D\neq0$ satisfies
  \[
   \mathrm{Lop}\ \sigma\ A\ H\ x\ D\ B\ \le\ 0 ,
  \]
  with $\mathrm{hgrad}$ as in \noderef{hgrad} and $\mathrm{Lop}$ as in \noderef{Lop}.
\end{enumerate}
Here $x\in\mathrm{argmax}_\Omega(u-\varphi)$ means, as in (2.8), that $x\in\Omega$
and $(u-\varphi)(y)\le(u-\varphi)(x)$ for every $y\in\Omega$.

This is Definition Definition~2.1 for the operator (1.8). The admissible test functions are
those of the class (2.11), $\mathcal A^{+}_x(u,\Omega)=\{\varphi\in C^2(\Omega):
x\in\mathrm{argmax}_\Omega(u-\varphi),\ \mathcal X\varphi(x)\neq0\}$, and the required
inequality is $\mathcal L\varphi(x)\le0$. The restriction
$\mathcal X\varphi(x)\neq0$ --- not merely $\nabla\varphi(x)\neq0$, which is weaker when
$m<n$ --- is what allows $A$ and $H$ to be undefined at the origin, and it makes the class
of sub-solutions larger than the one used by Bardi and Goffi.

Two points of translation. First, $\mathcal L\varphi(x)$ is written through the jet of
$\varphi$ at $x$: for $\Omega$ open, $\varphi\in C^2(\Omega)$ and $x\in\Omega$ a
second-order jet of $\varphi$ at $x$ exists and is unique, namely
$(D,B)=(\nabla\varphi(x),D^2\varphi(x))$, so quantifying over all jets of $\varphi$ at $x$
names that one pair, and then
$\mathrm{hgrad}\ \sigma\ x\ D=\mathcal X\varphi(x)$ and
$\mathrm{hhessS}\ \sigma\ x\ D\ B=\mathcal X\mathcal X^{s}\varphi(x)$, so the displayed
inequality is $\mathcal L\varphi(x)\le0$. Second, the operator (1.8) has no
zeroth-order term, so the value $\varphi(x)$ does not enter and no monotonicity in $r$ has
to be recorded. The paper also notes that, since $\varphi$ may be modified away from $x$,
requiring $x$ to be a global maximum point of $u-\varphi$ on $\Omega$ is equivalent to
requiring it to be a local one; the global form is used here, following (2.11). The notion is only ever used with
$\Omega$ an open set, as in Definition Definition~2.1, which is what makes the quantification
over second-order jets of $\varphi$ at $x$ non-vacuous. Writing
$\mathcal Lu\le0$ in $\Omega$ in the viscosity sense means that $u$ is a viscosity
sub-solution in this sense.
\end{definition}
